% !TEX root = ../../main.tex
% C5.6 — Thiết kế giao diện người dùng (RÚT GỌN 2026-09-29; bỏ nguyên tắc "ngôn ngữ nhất quán" vì trùng NFR-19; bản cũ ở report/archive/)
% Khớp frontend/src: constants/navigation.js (menu theo vai trò, HOME_BY_ROLE), routes/AppRoutes.jsx + guards.jsx,
% features/results/ResultViewer.jsx, CreateCaseDialog, AddToCaseDialog, features/search/ImageDropzone.jsx, SearchResults.jsx,
% components/person/PersonCrop.jsx, pages/viewer/OverviewPage.jsx; architect §9 (dùng được ở 390 px). Hình H12.
\section{Thiết kế giao diện người dùng}
\label{sec:5.6}

Giao diện được thiết kế trực tiếp trên ứng dụng web, không qua bước bản mẫu; Mục~\ref{sec:6.3} trình bày các màn hình chính.

\subsection{Cấu trúc màn hình theo vai trò}
\label{sec:5.6.1}

Mỗi vai trò vào trang chủ riêng và chỉ thấy màn hình của mình (Hình~\ref{fig:screen-flow}). Trong hình, mũi tên đậm chỉ trang mở sau khi đăng nhập, khung nét đứt là hộp thoại, còn các màn hình cùng vai trò nối với nhau qua menu bên trái.

\begin{figure}[!htbp]
\centering
\includegraphics[width=\textwidth]{H12-luong-man-hinh.png}
\caption{Sơ đồ luồng màn hình theo vai trò}
\label{fig:screen-flow}
\end{figure}

\begin{itemize}
    \item \textbf{Quản trị viên} có tám màn hình trong hai nhóm \emph{Quản trị} và \emph{Giám sát}.
    \item \textbf{Giám sát viên} làm việc chủ yếu trên \emph{Tìm kiếm người}, nơi biểu mẫu và kết quả cùng một trang để tìm lại nhiều lần; một kết quả mở hộp thoại chi tiết, từ đó lưu vào vụ việc. \emph{Vụ việc của tôi} gồm danh sách và chi tiết trên cùng một trang.
    \item \textbf{Quản lý} bắt đầu từ \emph{Tổng quan} và mở vụ việc ở \emph{Hồ sơ vụ việc}, là bản chỉ đọc của màn hình vụ việc của Giám sát viên.
\end{itemize}

\subsection{Các nguyên tắc thiết kế}
\label{sec:5.6.2}

Vì mọi kết quả đều cần người dùng xác nhận, giao diện trước hết hỗ trợ đánh giá bằng mắt (NFR-20): ảnh người đặt trong khung hình, người được tìm thấy được làm nổi bật khi có nhiều người, thứ hạng, camera và giờ xuất hiện được ưu tiên, còn điểm phù hợp chỉ là thông tin phụ.

Giao diện cũng giảm thao tác lặp: ảnh truy vấn đưa vào được nhiều cách, kể cả kéo một kết quả vào ô tìm kiếm (FR-O01), và kết quả sắp xếp, lọc lại được mà không cần truy vấn mới.

Khi không có kết quả, giao diện nêu lý do và gợi ý cách tiếp tục; lỗi dùng thông báo tiếng Việt của máy chủ, và khi dữ liệu đã đổi ở nơi khác thì người dùng được yêu cầu tải lại. Trên điện thoại rộng khoảng 390 điểm ảnh, menu trái thành danh sách chọn và bảng rộng được cuộn ngang.
